\documentclass[11pt,a4paper]{article}
\usepackage[T1]{fontenc}
\usepackage{lmodern,amsmath,amssymb,microtype,booktabs}
\usepackage[margin=25mm]{geometry}
\usepackage[colorlinks=true,urlcolor=blue,linkcolor=blue]{hyperref}
\newcommand{\dd}{\mathop{}\!\mathrm d}
\title{Beyond the right recoil saddle:\\a two-branch Barnes--Bessel diagnostic}
\author{Alan Day \and Chengheng\thanks{AI research assistant (OpenAI).}}
\date{Working research note --- 30 October 2026}
\begin{document}\maketitle
\begin{abstract}
Exact negative-side reflection and a compact-kernel limit lead to an
auxiliary contour with two modified-Bessel contributions. Finite-parameter
quadrature shows a large additional contribution for some imaginary
energy transfers. This revises the prospective global single-saddle
strategy and preserves the preceding local result. Neither a full
Barnes-to-Bessel limit nor a joint growing-mass theorem is proved here.
Numerical errors quoted are empirical diagnostics.
\end{abstract}
\section{Object and exact negative reflection}
The de Sitter radius is one. We retain the internal chiral graph amplitude
obtained in the preceding fixed-mass \(\kappa\to\infty\) limit, then study
its masses. No physical state, observable algebra, external insertion or
energy cycle is newly identified here. This is not a finite-\(G\) OTOC.
For \(u=v=0\), \(K=\nu^2/(2M)>0\), \(\omega=iKy\), set
\[
 a_j=M+\nu/2\ (j=1,\ldots,4),\qquad
 r=(M+\nu+\omega/2,2M,M+\nu-\omega/2).
\]
The exact kernel and chiral integral are
\[
 B(z)=e^{-i\pi z}
 \frac{\prod_a\Gamma(z+ia)\prod_r\Gamma(\tfrac12-z-ir)}
 {\Gamma(z+\tfrac12)},\qquad
 \mathfrak B=\frac1{2\pi i}\int_{\mathbb R-ic}B(z)\,\dd z.
\]
For \(w=-z\), Euler reflection \cite{gamma} gives
\begin{align}
 B(-w)&=H_-(w)E_-(w),\\
 H_-(w)&=8\pi^3e^{4\pi iw+\pi\sum r}
 \frac{\Gamma(w+\tfrac12)\prod_r\Gamma(w+\tfrac12-ir)}
 {\prod_a\Gamma(1+w-ia)},\\
 E_-(w)&=\frac{1+e^{2\pi iw}}
 {\prod_a(1-e^{2\pi i(w-ia)})}.
\end{align}
No periodic denominator is discarded. This reflection identity also
holds for the general real transverse labels.
With \(f(t)=t(\log t-1)\), define
\[
 \Phi(z)=2\pi iz-\pi\sum r+\sum_a f(z+ia)-\sum_r f(z+ir)-f(z).
\]

\newpage
\section{Phase comparison and a computable contour}
For real \(a\), conjugation-invariant \(r\), and compatible
right-half-plane logarithms, the negative phase obeys
\[
 \Psi(w)=4\pi iw+\pi\sum r+f(w)+\sum_r f(w-ir)-\sum_a f(w-ia)
        =2\pi iw-\overline{\Phi(\bar w)}.
\]
Let \(z_+=x_+-ic_+\) be the right saddle. Write
\(\prod(s-a)-s\prod(s-r)=A_qs^2+bs+c_0\);
then \(s_+=(-b+i\sqrt{-\Delta})/(2A_q)\), \(z_+=-is_+\), \(\Delta<0\).
The difference of real phase values at conjugate points is
\[
 D=\Re\Phi(z_+)-\Re\Psi(\bar z_+)=2\pi c_++2\Re\Phi(z_+).
\]
This does not determine a contour coefficient.
For \(\epsilon=\nu/M\), \(\delta=M+\nu-c_+\), the symmetric
small-\(\epsilon\) expansion at bounded \(y\) is
\[
 D/K=2F(y)-2\pi\delta/K+O(\epsilon^2),\qquad
 F(y)=\sqrt{1+y^2}-y\,\operatorname{arsinh}y.
\]

Choose \(L=x_+>K|y|/2\), \(c=M+\nu-d\) with
\(\max a<c<\min\Re r\). The negative tail may be deformed downward
at \(x=-L\). The swept lower-left region has no poles:
the \(\Gamma(z+ia)\) poles lie at height \(-a>-c\), and the other
pole rows lie to the right of \(-L\).
Reflection and Stirling remove distant arcs at fixed parameters.
The orientation is upward on the left ray:
\[
 \int_{\mathbb R-ic}B(z)\,\dd z
 =\int_{-L}^\infty B(x-ic)\,\dd x
 +i\int_0^\infty B(-L-i(c+t))\,\dd t.
\]
This fixed-parameter deformation supplies no mass-uniform tail bound.

\section{Compact-kernel limit}
Put \(t=z+i(M+\nu)\), \(\alpha=Ky\), and
\[
 C_{M,\nu}=e^{-\pi(M+\nu)}
 \frac{\Gamma(-i\nu/2)^4\Gamma(\tfrac12-i(M-\nu))}
 {\Gamma(\tfrac12-i(M+\nu))}.
\]
At fixed \(K,\alpha\), \(M=\nu^2/(2K)\), \(\nu\to\infty\),
Gamma-ratio asymptotics \cite{stirling} give locally away from poles
\begin{align}
 \frac{B(t-i(M+\nu))}{C_{M,\nu}}&\longrightarrow U_{K,\alpha}(t),\\
 U_{K,\alpha}(t)&=e^{-2\pi it}(K/2)^{2t}
 \Gamma(\tfrac12-t+\alpha/2)\Gamma(\tfrac12-t-\alpha/2).
\end{align}
The leading powers are
\(( -i\nu/2)^{4t}[-i(M-\nu)]^{-t}[-i(M+\nu)]^{-t}e^{-i\pi t}\);
on these branches they tend to \(e^{-2\pi it}(K/2)^{2t}\).
This is compact convergence, not an exchanged full integral limit.

\newpage
\section{Exact auxiliary Bessel integral}
For real \(\alpha\), \(K>0\), choose \(\mathcal C\) from
\(d_0-i\infty\) to \(d_0+i\eta\), then right to \(+\infty+i\eta\),
where \(d_0<\tfrac12-|\alpha|/2\), \(\eta>0\).
Set \(J=(2\pi i)^{-1}\int_{\mathcal C}U\,\dd t\).
Its downward tail is a power times \(e^{-3\pi|\Im t|}\);
reflection on the right tail gives reciprocal Gamma factors and
factorial decay. Close clockwise below the horizontal segment,
with boundaries avoiding the real poles.

For noninteger \(\alpha\), summing the two pole families
\(t=\tfrac12\pm\alpha/2+n\) gives
\begin{align}
 J={}&-e^{-i\pi\alpha}\Gamma(-\alpha)(K/2)^{1+\alpha}
 {}_0F_1(1+\alpha;K^2/4)\nonumber\\
 &-e^{i\pi\alpha}\Gamma(\alpha)(K/2)^{1-\alpha}
 {}_0F_1(1-\alpha;K^2/4).
\end{align}
The residue sign includes both the derivative of the Gamma argument
and the clockwise closure. Using the \(I\) series and the standard
\(I_{-\alpha}-I_\alpha\) connection \cite{iseries,connection},
\begin{align}
 J&=\frac{\pi K}{2\sin(\pi\alpha)}
 \left(e^{-i\pi\alpha}I_\alpha(K)-e^{i\pi\alpha}I_{-\alpha}(K)\right)\\
 &=-K\left[e^{i\pi\alpha}K_\alpha(K)+i\pi I_\alpha(K)\right]
 =-K K_\alpha(Ke^{-i\pi}). \label{eq:J}
\end{align}
The last expression denotes continuation below the negative axis
\cite{continuation}. At integer order continue the \emph{complete sum};
the separate residue coefficients are singular.
This is an exact auxiliary identity, not yet the full Barnes limit,
and not a claim of a new special-function identity.

\subsection*{Two exponential scales}
For fixed \(y>0\), \(K\to\infty\), large-order expansions \cite{largeorder}
give
\[
 I_{Ky}(K)\sim\frac{e^{KF(y)}}{\sqrt{2\pi K\sqrt{1+y^2}}},
 \qquad
 K_{Ky}(K)\sim\sqrt{\frac{\pi}{2K\sqrt{1+y^2}}}e^{-KF(y)}.
\]
Thus \(K_{Ky}(K)/(\pi I_{Ky}(K))\sim e^{-2KF(y)}\).
Since \(F'(y)=-\operatorname{arsinh}y<0\), its positive zero is unique:
\[
 y_c\simeq1.5088795615383199289.
\]
This is an auxiliary exponential-magnitude crossover, not a proved
physical critical point or a Stokes transition of the full graph.
Phase cancellation of the complete sum must still be controlled.

\newpage
\section{A first correction and finite-parameter evidence}
At fixed \(K,\alpha,t\), the compact Gamma expansion gives
\[
 B/C_{M,\nu}=U_{K,\alpha}(t)
 \left[1+\frac{4i}{\nu}(t^2-t)+O_{K,\alpha,t}(\nu^{-2})\right].
\]
The four large arguments \(-i\nu/2\) supply
\(4t(t-1)/(2(-i\nu/2))=4i(t^2-t)/\nu\);
the remaining arguments and power corrections enter at order \(\nu^{-2}\).
Let \(\mathcal D=(K/2)\partial_K\), \emph{holding \(\alpha\) fixed}.
Then \(\mathcal DU=tU\), and the Bessel equation yields exactly
\[
 (\mathcal D^2-\mathcal D)J=(K^2+\alpha^2-1)J/4.
\]
Formal integration therefore predicts the normalized first-order phase
\(e^{i\alpha^2/\nu}\). This exponential completes a first-order term;
it is not a derived all-orders resummation or an integrated error theorem.

For \(u=v=0\), \(\nu=4K^3\), \(M=\nu^2/(2K)\), define
\[
 g_+=8\pi^3\prod_a(z_++ia)^{-1/2},\qquad
 \mathfrak B_{\rm loc,lead}=\frac{g_+e^{\Phi_+}}
 {i\sqrt{2\pi(-\Phi''_+)}}.
\]
The numerical full-contour to local-leading ratios are
\begin{center}\begin{tabular}{rrr}\toprule
\(K\)&\(y\)&\(\mathfrak B/\mathfrak B_{\rm loc,lead}\) (approximate)\\\midrule
4&0&\(1.0377303-0.0016775i\)\\
4&2&\(0.7574552-182.4665112i\)\\
8&0&\(1.0169041-0.0001656i\)\\
8&1&\(1.0014509-0.0003609i\)\\
8&2&\(-4.4614748-33541.8312194i\)\\
16&0&\(1.0081070-0.0000205i\)\\
16&2&\(-22852.0717-1.1232773362\,10^9i\)\\\bottomrule
\end{tabular}\end{center}
These finite values motivate a second contribution; they do not prove
an asymptotic counterexample along the entire sequence.
For \(R=\mathfrak B(K,Ky)/\mathfrak B(K,0)\),
\(R_J=J(K,Ky)/J(K,0)\), we find
\begin{center}\begin{tabular}{rrrr}\toprule
\(K\)&\(y\)&\(|R/R_J-1|\)&\(|R/(R_Je^{i(Ky)^2/\nu})-1|\)\\\midrule
4&2&0.24889358&0.00389292\\
8&1&0.03124995&0.00007430\\
8&2&0.12488755&0.00049349\\
16&2&0.06248789&0.00006189\\\bottomrule
\end{tabular}\end{center}
There is no fitted phase parameter. Its integrated validity is open.

\newpage
\section{Controls, limitations, and next steps}
We use 20/36 Gauss nodes per segment, 50/65 decimal digits, \(d=1\),
\(X=\max(32,8K)\), and downward extent \(T=18\).
For \(y=2\), additional runs use \(d=0.75\),
\(X=\max(32,10K)\), \(T=24\), and 90 digits.
The segment cancellation ratio reaches about \(60.5\);
all relative comparisons use the complete complex sum.
Extent changes test truncation empirically, without certifying
the omitted infinite tails.

All 34 checks pass: 12 reflection/phase identities,
6 auxiliary Bessel/differential identities, 7 quadrature refinements,
and 9 contour/extent/precision comparisons.
Parameters, segment values, refinement changes, Python and mpmath
versions, and code/input hashes are stored in
\path{results/stage4/recoil-global/diagnostics.json}.
The executable is \path{code/recoil_global_probe.py}.

The preceding local result is retained. The next objective is a
\emph{two-branch} analysis rather than a presupposed right-saddle
dominance for all bounded imaginary transfers.
First prove a full fixed-\(K,\alpha\) Barnes-to-Bessel limit.
The unresolved requirement is a mass-uniform integrable tail bound,
especially on the positive horizontal segment. A useful partition is
\(x\ll\nu\), \(\nu\lesssim x\lesssim M\), and \(x\gtrsim M\),
with overlap estimates and all exact periodic factors retained.
Compact convergence cannot replace this bound.

Only afterwards let \(K\) grow, specify its relation to \(\nu\),
and establish integrated correction estimates. Near comparable
branch magnitudes, control absolute errors before dividing by a
potentially small sum. Transverse \(u,v\) directions and exterior
tails require separate estimates.

The external energy cycles, insertion normalization, disconnected
baseline, and microscopic dictionary remain open.
Taking the fixed-mass \(\kappa\) limit before mass growth does not
establish a joint finite-\(G\) limit. Complex energies require analytic
continuation of the chiral factors, not a newly imposed modulus square.
This is an author self-reviewed research draft.

\begin{thebibliography}{9}
\bibitem{gamma} NIST DLMF, \S5.5, functional relations,
\url{https://dlmf.nist.gov/5.5}.
\bibitem{stirling} NIST DLMF, \S5.11, asymptotic expansions,
\url{https://dlmf.nist.gov/5.11}.
\bibitem{iseries} NIST DLMF, \S10.25, modified Bessel definitions,
\url{https://dlmf.nist.gov/10.25}.
\bibitem{connection} NIST DLMF, \S10.27, connection formulas,
\url{https://dlmf.nist.gov/10.27}.
\bibitem{continuation} NIST DLMF, \S10.34, especially 10.34.2,
\url{https://dlmf.nist.gov/10.34}.
\bibitem{largeorder} NIST DLMF, \S10.41(ii), large-order expansions,
\url{https://dlmf.nist.gov/10.41}.
\end{thebibliography}
\end{document}
